\chapter{与经典上同调的比较：光滑概形的情形}
%\addtocontents{toc}{M. Artin}
\begin{center}
M. Artin
\end{center}
\phantomsection\MakeLinkTarget*{sga4.exp.XI}\label{sga4.exp.XI}

\setcounter{pageoriginal}{63}

\section{引言}\etiquette[1]{XI.sec1}\pageoriginale

在第~\hyperref[XI.sec4]{4} 节中，我们将证明：对于 $\Spec\CC$ 上的光滑概形 $X$，以局部常值挠层（locally constant torsion sheaf）为系数的平展上同调（étale cohomology）与经典上同调（classical cohomology）一致。这个证明本质上就是 \cite{sga4.exp.XI.bib.1} 中给出的证明；只要有下面回顾的 Grauert 定理（\hyperref[XI.thm4.3]{4.3} (iii)），即 $X(\CC)$ 的每个有限覆盖都来自 $X$ 的一个平展覆盖，证明就是初等的。因此，尽管稍后我们将利用 Hironaka 的奇点消解定理 \cite{sga4.exp.XI.bib.3} 证明一个更完整的结果，这里仍给出这个证明\footnote[1]{还应指出，Hironaka 的奇点消解也给出了 Grauert 结果的一个非常简单的证明；从这个意义上说，无论如何都可以把它看作比较定理的基础。关于这一点，参见 Mme M. Raynaud 在 \SGA 1 XII（North Holland Pub. Cie）中的报告。}。此外，第~\hyperref[XI.sec3]{3} 节中构造的良好邻域，在其他\pageoriginale 情形中也可能有用。

注意，对于非挠层，平展上同调与经典上同调并不一致。例如，对每个正规概形 $X$ 都有 $H^1(X_{\et},\ZZ)=0$（\hyperrefext[IX][IX.prop3.6]{3.6} (ii)），但经典上同调未必为零。

\section{充分一般的超平面截面的存在性}\etiquette[2]{XI.sec2}

下面的结果是经典的，但尚缺少文献出处（有待 \EGA V）：

\begin{enonce*}{定理 2.1}\etiquette[2.1]{XI.thm2.1} %Why is the statement
                                %not formulated for a linear subspace
                                %varying around a fixed point?
                                %That is, (i) with $P\in L$ or
                                %$(ii)$ with $d=1$
设 $k$ 是一个域，$V\subset \PP^n_k$ 是一个纯 $r$ 维的射影代数概形。设 $V'$ 是由 $V$ 中在 $\Spec k$ 上光滑的点组成的开子概形，$P$ 是 $\PP^n$ 的一个有理点。则：
\begin{enumerate}
\renewcommand{\theenumi}{\roman{enumi}}
\renewcommand{\labelenumi}{(\theenumi)}
\item $\PP^n$ 中给定维数 $d$ 的一个“充分一般”的线性子空间 $L$，在 $L\cap V'$ 的每一点都与 $V$ 横截相交（meet transversely）。

\item 设 $\{H_1,\ldots,H_s\}$ $(s\leq r)$ 是一组 %a
                                %'hyperplane' is not linear here
  经过 $P$ 的 $d\geq 2$ 次超曲面，并设 $Y=H_1\cap\ldots \cap H_s$。若 $\{H_1,\ldots,H_s\}$ 充分一般，则 $Y$ 在 $Y\cap V'$ 的每一点都与 $V$ 横截相交。
\end{enumerate}
\end{enonce*}

“对于充分一般的线性子空间”是指：存在参数化 $d$ 维线性子空间的 Grassmann 概形的一个稠密开子集 $U$，只要 $L$ 的参数属于 $U$，所讨论的关于 $L$ 的断言就成立。因此，当 $k$ 不是无限域时，这并不蕴含这样的（在 $k$ 上定义的）$L$ 存在。同样，
% mention the Veronese embedding?
经过 $P$ 的 $d$ 次超曲面由一个射影空间 $\PP^N$ \pageoriginale 参数化，而“对于充分一般的 $\{H_1,\ldots,H_s\}$”是指：存在 $(\PP^N)^s$ 的一个稠密开子集 $U$，只要 $(H_1,\ldots,H_s)$ 的参数属于 $U$，所讨论的断言就成立。

(i) 的证明：由归纳法，立即归结为 $L$ 是超平面的情形。设 $(V'_{\nu})$ 是 $V'$ 的一个有限开覆盖。如果对每个 $\nu$，将符号 $V'$ 换成 $V'_{\nu}$ 后断言 (i) 都成立，那么它对 $V'$ 也成立，因为有限个稠密开子集的交仍是稠密开子集。因此，我们可以把问题换成仿射代数概形的情形，记作 $V\subset \EE^n_k=\Spec k[x_1,\ldots,x_n]$。超平面 $L$ 将是一个线性方程的零点集：
$$
\ell(x)=a_0+\sum^n_{i=1}a_ix_i{\quoi}.
$$

回顾光滑概形局部是完全交。因此，我们可以把 $V'$ 换成 $V$ 中点的一个开子集，使得存在在 $V$ 上为零的多项式 $f_1(x),\ldots,f_{n-r}(x)$，以及指标 $1\leq j_{\nu}\leq n$ $(\nu=1,\ldots,n-r)$，使行列式
$$
\left|\frac{\partial f_i}{\partial x_{j_{\nu}}}\right|\quad 
(i,\nu=1,\ldots,n-r)
$$
可逆。不妨设 $j_{\nu}=\nu$ $(\nu=1,\ldots,n-r)$。于是，我们可以由方程
$$
y_j=\sum^{n-r}_{i=1}\quad c_i\quad \frac{\partial f_i}{\partial x_j}
$$
解出\pageoriginale $c_i$，将其表示为变量 $x_i$、$y_j$ $(i=1,\ldots, n-r; j=1,\ldots,n)$ 的有理函数，而且这些函数 $c_i(x,y)$ 在 $V'$ 的每一点都有定义。

现在，$L$ 在 $V'$ 的一点 $Q$ 与 $V$ 横截相交，等价于说矩阵
$$
\begin{bmatrix}
\frac{\partial f_1}{\partial x_1} & \cdots & \frac{\partial
  f_1}{\partial x_n}\\
\vdots & &\vdots\\
\frac{\partial f_{n-r}}{\partial x_1} & & \frac{\partial
  f_{n-r}}{\partial x_n}\\
a_1 & \cdots & a_n
\end{bmatrix}
$$
在 $Q$ 的值的秩为 $n-r+1$。所以，$L$ 在 $V'$ 的每一点与 $V$ 横截相交的充要条件，是由以下 $r+1$ 个方程组成的方程组
\begin{align*}
a_j &= \sum^{n-r}_{i=1} c_i (x,a)\frac{\partial f_i}{\partial x_j}\qquad
(j=n-r+1,\ldots, n)\\
a_0 &+ \sum^n_{i=1} a_i x_i =0
\end{align*}
在 $V'$ 上无解。这样的条件显然关于 $(a)$ 是可构造的：设 $X=V'\times A$，其中 $A$ 是以 $a_0,\ldots,a_n$ 为坐标的仿射空间，并设 $Y\subset X$ 是这些方程的解组成的闭子集。所述条件就是 $Y/A$ 的纤维为空，这确实是一个可构造\pageoriginale 条件（\EGA IV 9.5.1）。因此，只需 % \ref
验证 $A$ 的一般点上的纤维为空，而这由方程的形式立即得到。(i) 得证。(ii)：考虑由经过 $P$ 的 $d$ 次（$d\geq 2$）超曲面线性系给出的有理映射 $\PP^n\rightarrow \PP^N$。众所周知（也容易验证），这是 $P$ 的一个“爆破”（blowing up），因此给出一个局部闭浸入
$$
f:\PP^n-\{P\}\longrightarrow \PP^N{\quoi}.
$$
设 $V''=V'-(\{P\}\cap V')$，$\bar{V}{}''=\sisi{f}{f}(V'')$，并设 $L_i$ 是 $\PP^N$ 中对应于超曲面 $H_i$ 的超平面 $(i=1,\ldots,s)$。则 $V''$ 与 $\bar{V}{}''$ 同构，而且根据 $f$ 的定义，$Y$ 与 $V''$ 横截相交，等价于 $L_1\cap\ldots \cap L_s$ 与 $\bar{V}{}''$ 横截相交。令 $\bar{V}=$ $\bar{V}{}''$ 的闭包，便归结为对 $\bar{V}$ 应用断言 (i)，由此得到将 $V'$ 换成 $V''$ 后的 (ii)。还需考虑点 $P$（如果它是 $V$ 的光滑点）。但显然，若 $\{H_1,\ldots,H_s\}$ 充分一般且 $P$ 是 $V$ 的光滑点，则 $Y$ 在 $P$ 与 $V$ 横截相交，结论得证。


\section{良好邻域的构造}\etiquette[3]{XI.sec3}

\begin{enonce*}{定义 3.1}\etiquette[3.1]{XI.def3.1}
初等纤维化（elementary fibration）是一个概形态射 ${f:X\rightarrow S}$，它可以嵌入一个交换图
$$
\xymatrix@=2cm{
X  \ar[r]^{j} \ar[rd]_{f}& \bar{X} \ar[d]^{\bar{f}}& Y \ar[l]_{i}  \ar[ld]^{g}\\
               &    S   &      
}
$$
并满足\pageoriginale 下列条件：
\begin{enumerate}
\renewcommand{\theenumi}{\roman{enumi}}
\renewcommand{\labelenumi}{(\theenumi)} 
% \item $i$ is a closed immersion.
\item $j$ 是开浸入，在每个纤维中稠密，且 $X=\bar{X}-Y$。

\item $\bar{f}$ 是光滑射影态射，其纤维为几何不可约的 $1$ 维概形。

\item $g$ 是平展覆盖，而且 $g$ 的每个纤维均非空。
\end{enumerate}
\end{enonce*}


容易验证，$X$ 的这样的嵌入如果存在，则在典范同构意义下唯一。

\begin{enonce*}{定义 3.2}\etiquette[3.2]{XI.def3.2}
相对于 $S$ 的良好邻域（good neighborhood）是一个 $S$-概形 $X$，使得存在 $S$-概形
$$
X=X_n,\ldots,X_0=S
$$
以及初等纤维化 $f_i:X_i\rightarrow X_{i-1}$，
% $f_i$ is an $S$-morphism
$i=1,\ldots, n$。
\end{enonce*}

\begin{enonce*}{命题 3.3}\etiquette[3.3]{XI.prop3.3}
\footnote{这个结果容易推广为半局部的断言。}
设 $k$ 是代数闭域，$X/\Spec k$ 是光滑概形，$x\in X$ 是一个有理点。则存在 $X$ 的一个包含 $x$ 的开子集，它是一个良好邻域（相对于 $\Spec k$）。
\end{enonce*}

证明。取 $X$ 不可约。对 $\dim X=n$ 作归纳，只需找到 $x$ 在 $X$ 中的一个邻域 $U$ 和一个初等纤维化 $f:U\rightarrow V$，使得 $V$ 光滑且维数为 $n-1$。事实上，将存在 $v=f(x)$ 的一个良好邻域 $V'$，于是可取 $U'=U\cap \sisi{f^{-1}(V')}{V'}$ 作为 $x$ 的良好邻域。

可以假定 $X\subset \EE^r$ 是仿射的。设 $X_0$ 是 $X$ 在 $\PP^r$ 中的闭包。设 $\bar{X}$ 是 $X_0$ 的正规化，并设 $Y=\bar{X}-X$，赋予它诱导的约化结构。此外，设\pageoriginale $S\subset \bar{X}$ 为奇异点组成的闭子集。有 $S\subset Y$，且
\begin{align*}
\dim \bar{X} &= \dim X=n{\quoi},\\
\dim Y &= n-1{\quoi},\\
\dim S &\leq n-2{\quoi}.
\end{align*}

利用形如 $M^{\otimes r}$ 的可逆层 $L$ 将 $\bar{X}$ 嵌入射影空间 $\PP^N$，其中 $M$ 很丰沛，$r\geq 2$。于是存在 $\PP^N$ 中的超平面 % linear
$H_1,\ldots, H_{n-1}$，其中 $H_i$ 是方程
$$
\sum^N_{\nu =0}a_{i\nu} x_{\nu} =0{\quoi},
$$
的零点集，使它们包含 $x$，而且交集 $L=H_1\cap\ldots \cap H_{n-1}$ 的维数为 $N-n+1$，并与 $\bar{X}$、$Y$ 横截相交（\hyperref[XI.thm2.1]{2.1}）。
% see the comment on Theorem 2.1
 %For the intersection with $\bar{X}$: $S\cap L=∅$ for sufficiently
 %general $L$, hence $\bar{X}\cap L=(\bar{X} - Y)\cap L$ and
 %$\bar{X}- Y$ is smooth.
 %For the intersection with $Y$: $Y$ is reduced and $k$ algebraically closed,
 %hence every irreducible component is generically smooth, and sufficiently
 %general $L$ does not meet the singular locus of $Y$.
交集 $\bar{X}\cap L$ 是一条光滑连通曲线（这由 Bertini 定理得到），而 $Y\cap L$ 的维数为 $0$。再选取另一个超平面
$$
H_0:\sum^N_{\nu=0}a_{0 \nu}x_{\nu}=0
$$
使得 $H_0$ 与 $\sisi{\bar{X}}{X}\cap L$ 横截相交，且 $H_0\cap Y\cap L=∅$。 % One must also choose $x\notin H_0$ so that
                    % the map $f'$ below is well defined at $x$

考虑通过“引入射影坐标”
$$
y_i=\sum^N_{\nu =0}a_{i_{\nu}}x_{\nu}{\quoi}.
$$
得到的投影 $\PP^N\rightarrow \PP^{n-1}$。这是在投影中心 $C=H_0\cap\ldots \cap H_{n-1}$ 之外有定义的有理映射。设 $\epsilon :P'\rightarrow \PP^N$ 为沿 $C$ 的爆破，从而有图
$$
\xymatrix@=1.5cm{
\PP^{N}  \ar[rd]& &P'\ar[ld]^{\pi}\ar[ll]_{\epsilon}\\
               &\PP^{n-1} &
}
$$
其中\pageoriginale $\pi$ 为一个态射。设 $\bar{X}{}'\subset P'$ 为 $\bar{X}$ 的“严格”逆像（proper inverse image），即 $\epsilon^{-1}(\bar{X}-(\bar{X}\cap C))$ 的闭包。由于按假设 $C$ 与 $\bar{X}$ 横截相交，态射 $\bar{X}{}'\rightarrow \bar{X}$ 是沿有限集 $\bar{X}\cap C$ 的爆破。设 $X'=X-\bar{X}\cap C$，也将它等同于 $\bar{X}{}'$ 的一个开子概形；并设 $Y'=\bar{X}{}'-X'$ 为赋予诱导约化结构的闭子概形。有如下态射图 $\DD$
$$
\xymatrix@=2cm{
  X'  \ar[r]^{j} \ar[rd]_{f'}& \bar{X}{}'  \ar[d]^{\bar{f}{}'} & Y'\ar[l]_{i} \ar[ld]^{g'}     &\\
                           & \PP^{n-1} &\qquad ,   
}
$$
而我断言，存在 $v=f'(x)$ 的一个邻域 $V$，使 $\DD$ 限制到 $V$ 后满足（\hyperref[XI.def3.1]{3.1}）中的条件 (i)、(ii)、(iii)，因而 $f'_{\mid V}$ 是一个初等纤维化。这将完成证明。

条件 (i) 是平凡的。对于 (ii)，注意按假设 $\bar{X}\cap L$ 是一条光滑曲线，而且立即可见，$\epsilon$ 诱导一个一一态射 ${\bar{f}{}'}^{-1}(v)\rightarrow \bar{X}\cap L$。
% one-to-one morphism?
因此，$\left({\bar{f}{}'}^{-1}(v)\right)_{\red}\rightarrow \bar{X}\cap L$ 是双射。为了验证 $\bar{f}{}'$ 在 $v$ 的一个邻域上光滑，
% remove the \left and \right?
由 Hironaka 引理（\SGA 1 II 2.6 %\ref
\footnote{以及 \EGA IV 5.12.10。}），只需验证它在 ${\bar{f}{}'}^{-1}(v)$ 的一般点光滑。在这一点，$\bar{X}{}'$ 与 $\bar{X}$ 同构，而该态射光滑，因为 $L$ 与 $\bar{X}$ 横截相交。

还需证明 $g'$ 在 $v$ 的一个邻域上是平展的（因为 $Y$ 的维数为 $n-1$，显然 $g'$ 的每个纤维均非空）。有 $Y'=\epsilon^{-1}(Y)\amalg D_1\amalg\ldots \amalg D_r$，其中 $\bar{X}\cap C=\{P_1,\ldots,P_r\}$，而 $D_i$ 是 $\bar{X}$ 中 $P_i$ 的爆破\pageoriginale。每个 $D_i$ 等同于 $\epsilon^{-1}(P_i)$，且同构地映到 $\PP^{n-1}$ 上\sisi{。此外，投影}{}在 $Y$ 的每一点都有定义%since $Y\cap C=∅$
，而且在 $Y$ 上诱导的态射就是 $g'_{\mid \epsilon^{-1}(Y)}$。由于 $L$ 与 $Y$ 横截相交，它在 $v$ 上是平展的，因此 $g'$ 在~$v$ 上平展。


\section{比较定理}\etiquette[4]{XI.sec4}

\setcounter{subsection}{-1}
\subsection{}\etiquette[4.0]{XI.subsec4.0}
设 $X$ 是 $\Spec \CC$ 上局部有限型的概形，并在 $X$ 的有理点空间 $X(\CC)$ 上考虑通常拓扑。以 $X_{\cl}$ 表示由\emph{局部同构}（local isomorphism）$f:U\rightarrow X(\CC)$ 组成的景，即由满足下述性质的拓扑空间态射 $f$ 组成的景：

\sisi{对}{对}每个 $x\in U$，存在开邻域 $U_x$，使得 $f_{\mid U_x}$ 是 $U_x$ 到 $f(x)$ 的一个开邻域的同胚。

$X_{\cl}$ 中的一族态射 $\{U_{\nu}\rightarrow U\}$ 称为\emph{覆盖}的，如果 $U$ 是各 $U_{\nu}$ 的像的并。

由于开浸入是局部同构，我们有一个景的态射（包含）（IV\quad ） % a reference is missing
\footnote{更确切地说，范畴的包含 $\Ouv(X(\CC))\rightarrow X_{\cl}$ 定义了反方向的景态射 \hyperref[XI.eq4.1]{(4.1)}。}
\begin{equation*}
\delta:X_{\cl}\longrightarrow X(\CC)\qquad
(\ast){\quoi}.\tag{4.1}
\end{equation*}\etiquette[4.1]{XI.eq4.1}
现在，对于每个局部同构 $U\rightarrow X(\CC)$，按定义，都存在由 $X(\CC)$ 的开子集给出的 $U$ 的一个“覆盖”，由（\quad ）可推出，这两个景所关联的拓扑斯 % a reference is missing
是等价的（通过 $\delta_{\ast}$）。特别地，计算通常上同调时，可以把 $X(\CC)$ 换成 $X_{\cl}$。

再设 $f:X'\rightarrow X$ 是一个平展概形态射。则 $f(\CC):X'(\CC)\rightarrow X(\CC)$ 是局部同构。这由\pageoriginale Jacobian 判别法（\SGA 1 II 4.10）和隐函数定理立即得到。因此，函子 $X'\mapsto X'(\CC)$ 给出一个景的态射
\begin{equation*}
\epsilon:X_{\cl}\longrightarrow X_{\et}{\quoi},\tag{4.2}
\end{equation*}\etiquette[4.2]{XI.eq4.2}
其中 $X_{\et}$ 是平展景。

回顾以下结果\footnote{这些结果见第~2 页脚注所引的那一讲。}： % identify the Exposés
                                % above to identify the Exposé
                                % mentioned in the note?

\begin{enonce*}{定理 4.3}\etiquette[4.3]{XI.thm4.3}
\begin{enumerate}
\renewcommand{\theenumi}{\roman{enumi}}
\renewcommand{\labelenumi}{(\theenumi)}
\item $X$ 连通且非空，当且仅当 $X(\CC)$ 连通且非空。

\item 函子 $\epsilon$ 是全忠实的。

\item （“Grauert--Remmert 定理”）。函子 $\epsilon$ 在下述两个范畴之间诱导等价：赋予局部紧空间的通常拓扑的 $X(\CC)$ 的有限覆盖范畴，以及 $X$ 的平展覆盖范畴。
\end{enumerate}
\end{enonce*}

我们把 (i) 作为已知结果接受。例如，它可由 (GAGA)
%\ref
\cite{sga4.exp.XI.bib.4} 推出。断言 (ii) 容易由 (i) 得到。事实上，为了验证 (ii) 中包含的满射性断言，设 $\varphi:X'(\CC)\rightarrow X''(\CC)$ 是一个 $X(\CC)$-态射。可以假定所有概形都是连通仿射的，因而是分离的。于是，$\varphi$ 的图像 $\Gamma$ 是 $\fprod{X'}{X''}{X}(\CC)$ 的一个连通分支，由 (i)，存在 $\fprod{X'}{X''}{X}$ 的某个连通分支 $Y$，使得 $\Gamma=Y(\CC)$。于是 $Y$ 是诱导 $\varphi$ 的一个态射 $f:X'\rightarrow X''$ 的图像，即投影 $Y\rightarrow X'$ 是同构：事实上，$Y(\CC)\rightarrow X'(\CC)$ 是同构，由此容易推出 $Y\rightarrow X'$ 也是同构（它是平展、根性（radicial）且满的！）。对于 (iii)，由 (ii)，问题完全归结为证明 $X(\CC)$ 的每个有限覆盖都形如 $X'(\CC)$，其中 $X'\rightarrow X$\pageoriginale 是平展的。

为此，利用下降（\SGA 1 IX 4.7）立即归结为
%\ref
$X$ 连通且正规的情形。该问题关于 $X$ 是局部的（对 Zariski 拓扑而言），因此可以假定 $X$ 是仿射的。设 $f:X\rightarrow \EE^n$ 是有限态射，$T'\rightarrow X(\CC)$ 是连通有限覆盖。立即可验证，复合态射 $\eta':T'\rightarrow \EE^n(\CC)$ 是（\cite{sga4.exp.XI.bib.2} \S 2，定义~3）意义下的“解析覆盖”（analytic covering）。由（\cite{sga4.exp.XI.bib.2} \S 2，Satz 8），\sisi{$\eta'$}{}可扩张为一个“解析覆盖” $\eta:T\rightarrow \PP^n(\CC)$，而只需证明 $\eta$ 是由适当的概形态射 $f:Y\rightarrow \PP^n$ \sisi{诱导}{诱导}的，其中 $Y$ 正规。现在，基本定理（\cite{sga4.exp.XI.bib.2} \S 13 Satz 42）断言，$T$ 上存在一个典范的正规解析结构，使得 $\eta$ 诱导一个解析空间态射。由于此时 $\eta$ 是有限态射，$T$ 上结构层的正像是 $\PP^n$ 上的一个凝聚解析代数层 $\mathfrak{a}$。由 \sisi{(GAGA)}{GAGA} \cite{sga4.exp.XI.bib.4} % \ref
可知，$\mathfrak{a}$ 是与 $\PP^n$ 上某个代数层 $A$ 关联的解析层，于是取 $Y=\sheaf{\Spec}A$。

现在可以证明以下结果：

\begin{enonce*}{定理 4.4}\etiquette[4.4]{XI.thm4.4}
设 $X$ 是 $\Spec \CC$ 上的光滑概形。
\begin{enumerate}
\renewcommand{\theenumi}{\roman{enumi}}
\renewcommand{\labelenumi}{(\theenumi)}
\item $X_{\et}$ 上的可构造局部常值挠层范畴，与 $X_{\cl}$ 上具有有限纤维的局部常值挠层范畴之间存在等价，此等价由互为拟逆的函子 $\epsilon_{\ast}$ 和 $\epsilon^{\ast}$ 给出。

\item 设 $F$ 是 $X_{\cl}$ 上具有有限纤维的局部常值挠层，并仍以 $F$ 表示 $X_{\et}$ 上诱导的层 $\epsilon_{\ast}F$。则
$$
R^q\epsilon_{\ast}F=0\quad \mbox{ 若 } q>0{\quoi}.
$$

\item 使用\pageoriginale {\rm(i)} 中的记号，典范态射
$$
H^q(X_{\et},F)\longrightarrow H^q(X_{\cl},F)
$$
对每个 $q$ 都是双射。特别地，对每个 $q$ 都有
$$
H^q(X_{\et}, \ZZ/\sisi{n\ZZ}{n})\xrightarrow{\sim} H^q(X_{\cl},\ZZ/\sisi{n\ZZ}{n})
$$
成立。
\end{enumerate}
\end{enonce*}

证明。对于 (i)，注意所讨论的层就是能由 $X$ 的平展覆盖（\resp $X(\CC)$ 的有限覆盖）表出的层。由于这些覆盖的范畴之间存在等价（\hyperref[XI.thm4.3]{4.3}），断言 (i) 立即得到。(iii) 中的同构由 (ii) 及 $\epsilon$ 的 Leray 谱序列 V \hyperref[V.prop5.2]{5.2} % \ref + 5.3 instead of 5.2?
得到。还需证明 (ii)。

现在，$R^q\epsilon_{\ast}F$ 是与预层 $\Rr^qF,$ 关联的层，其中，对平展态射 $X'\rightarrow X$，有 $(\Rr^qF)(X')=H^q(X'_{\cl},F)$。因此，为了证明当 $q>0$ 时 $R^q\epsilon_{\ast}F=0$，必须证明以下\sisi{引理}{}：

\begin{enonce*}{引理 4.5}\etiquette[4.5]{XI.lem4.5}
设 $\xi\in H^q(X_{\cl},F)$，$x\in X_{\cl}=X(\CC)$。存在像包含 $x$ 的平展态射 $X'\rightarrow X$，使得 $\xi$ 在 $H^q(X'_{\cl},F)$ 中的像为零。
\end{enonce*}

引理的证明。对 $n=\dim X$ 作归纳；问题关于 $X$ 对平展拓扑是局部的，\sisi{}{因而}可以假定 $F$ 为常值层，再利用分步化简（dévissage），甚至可以假定 $F=\ZZ/\sisi{n\ZZ}{n}$。此外，将 $X$ 换成 $x$ 的一个 Zariski 开邻域后，可以假定（\hyperref[XI.prop3.3]{3.3}）$X$ 有一个初等纤维化（\hyperref[XI.def3.1]{3.1}）$f:X\rightarrow S$。沿用（\hyperref[XI.def3.1]{3.1}）的记号：通过\pageoriginale 直接计算，可见当 $q>1$ 时有 $R^q{j}_{\ast}F=0$，$R^1j_{\cl\!\ast}F$ 是 $Y_{\cl}$ 上一个常值层的补零延拓（extension by $0$），最后，$j_{\cl\!\ast}F$ 是 $\bar{X}$ 上的常值层 $\ZZ/\sisi{n\ZZ}{n}$。此外，可以逐纤维计算 $R^q \bar{f}_{\cl\!\ast}(j_{\cl\!\ast}F)$。由 Leray 谱序列
$$
E^{pq}_2=R^p\bar{f}_{\cl\!\ast}(R^qj_{\cl\!\ast}F)\Longrightarrow
R^{p+q}f_{\cl\!\ast}F
$$
可推出 $R^{p+q}f_{\cl\!\ast}F$ 同样可以逐纤维计算，因此

$R^0f_{\cl\!\ast}\ZZ/\sisi{n\ZZ}{n}=\ZZ/\sisi{n\ZZ}{n}$，

$R^1f_{\cl\!\ast}\ZZ/\sisi{n\ZZ}{n}$ 是 $S_{\cl}$ 上的一个局部常值挠层，

当 $q>1$ 时，$R^qf_{\cl\!\ast}\ZZ/\sisi{n\ZZ}{n}=0$。

于是，Leray 谱序列

$E^{pq}_2=H^p(S_{\cl}, R^qf_{\cl\!\ast}F)\Longrightarrow
H^{p+q}(X_{\cl},F)$

\noindent
化为正合列
\begin{equation*}
\ldots\longrightarrow H^q(S_{\cl},f_{\cl\!\ast}F)\longrightarrow
H^q(X_{\cl},F)\longrightarrow H^{q-1}(S_{\cl},
R^1f_{\cl\!\ast}F)\longrightarrow\tag{$\ast$}
\end{equation*}
设 $s$ 为 $x$ 在 $S$ 中的像。由归纳假设及 (i)，对于每个类 $\eta\in H^q(S_{\cl},G)$（$G$ 为具有有限纤维的局部常值挠层），存在 $s$ 的一个“平展邻域” $S'\rightarrow S$，使 $\eta$ 在 $H^q(S'_{\cl},G)$ 中的像为零。现在，函子 $R^qf_{\cl\!\ast}$ 显然与平展基变换可交换，而且若 $S'\rightarrow S$ 是 $s$ 的平展邻域，则 $X'=\fprod{X}{S'}{S}\rightarrow X$ 是 $x$ 的平展邻域。因此，当 $q\geq 1$ 时，利用正合列 $(\ast)$，先取 $G=R^1f_{\cl\!\ast}F$，再\pageoriginale 取 $G=f_{\cl\!\ast}F$，便得到结论。 % The proof ends here; the
                                % following is a remark for the case $q=1$
当 $q=1$ 时，注意 $H^1(X_{\cl},\ZZ/n \ZZ)$（\resp $H^1(X,\ZZ/n\ZZ)$）分类 $X_{\cl}$（\resp $X$）的群为 $\ZZ/n\ZZ$ 的主平展覆盖；
% should one instead speak of finite coverings for $X_\cl$?
因而由 \hyperref[XI.thm4.3]{4.3} (iii)，映射 $H^1(X,\ZZ/n\ZZ)\rightarrow H^1(X_{\cl},\ZZ/n\ZZ)$ 是双射，由此立即得到 $q=1$ 时的消去结果 \hyperref[XI.lem4.5]{4.5}。

\begin{enonce*}[remark]{变体 4.6}\etiquette[4.6]{XI.var4.6}
还可以用下面这个归于 Serre 的巧妙方法证明（\hyperref[XI.lem4.5]{4.5}）：注意一个连通良好邻域（\hyperref[XI.def3.2]{3.2}）$X$ 满足下列两个条件：
\begin{enumerate}
\renewcommand{\theenumi}{\roman{enumi}}
\renewcommand{\labelenumi}{(\theenumi)}
\item 若 $n>1$，则 $\pi_n(X(\CC))=0$。

\item $\pi_1(X(\CC))$ 是有限型自由群的迭次扩张。
\end{enumerate}
\end{enonce*}

事实上，若 $X\rightarrow S$ 是一个初等纤维化，其中 $S$ 是良好邻域，则由归纳法，可以假定 (i)、(ii) 对 $S$ 成立。现在，容易看出 $X(\CC)\rightarrow S(\CC)$ 是局部平凡纤维空间，其纤维同构于 $X_0(\CC)$，其中 $X_0$ 为 $X/S$ 的任意一个纤维。此时 $X_0$ 是连通非奇异的\emph{非完备}曲线，而众所周知，这蕴含当 $n>1$ 时 $\pi_n(X_0(\CC))=0$，且 $\pi_1(X_0(\CC))$ 是自由群。 % \ref
于是 (i)、(ii) 可由同伦正合列得到：
$$
\longrightarrow \pi_n(X_0(\CC))\longrightarrow
\pi_n (X(\CC))\longrightarrow \pi_n
(S(\CC))\longrightarrow\ldots
$$

所以 $X(\CC)$ 是一个 $K(\pi,1)$ 空间，这蕴含每个上同调类 $\xi\in H^n(X(\CC),F)$ 在 $X(\CC)$ 的万有覆盖上变为零。只需证明 $\xi$ 在 $X$ 的某个有限覆盖 $X'$ 上已经变为零，这归结为证明 $\pi_1=\pi_1(X(\CC))$ 是\pageoriginale 一个\emph{好群}（good group）（\cf CG 1 16）， %\ref
也就是说，从 $\pi_1$ 到其完备化 $\sisi{\widehat}{}{\pi}_1=\varprojlim_N \pi_1/N$ 的态射（其中 $N$ 遍历 $\pi_1$ 的所有有限指数不变子群）
% 'invariant' = 'normal'?
对每个连续有限 $\pi_1$-模 $M$ 都诱导一个双射
$$
H^q(\widehat{\pi}_1,M)\xrightarrow{\sim} H^q(\pi_1,M)
$$
（其中左边的符号表示 $\widehat{\pi}_1$ 作为概有限群的上同调）。而有限型自由群的迭次扩张是好群（\cf CG I p.~15-16, exc.~1,2），所以结果得证。 % \ref

\begin{thebibliography}{99}
\bibitem{sga4.exp.XI.bib.1} GIRAUD J., ``Analysis Situs'', Séminaire
Bourbaki, Exposé nᵒ~256.

\bibitem{sga4.exp.XI.bib.2} GRAUERT H\sisi{.}{} and REMMERT R., ``Komplexe Räume'',
Math, Ann. Bd. 136, p. 345 (1958).

\bibitem{sga4.exp.XI.bib.3} HIRONAKA H., ``Resolution of singularities of an
algebraic variety over a field of characteristic zero'' Annals of
Math. Vol. 79, nᵒ~1, 2 (1964).

\bibitem{sga4.exp.XI.bib.4} SERRE J.P., ``Géométrie Algébrique et Géométrie
Analytique'', Annales de l'Institut Fourier, t. VI, p. 1-12 (1956).
\end{thebibliography}
